\uuid{TiZs}
\titre{Approximation normale d'une somme de durées exponentielles}
\auteur{Maxime Nguyen}
\organisation{AMSCC}
\datecreate{2026-11-27}
\niveau{L2}
\module{Probabilité}
\chapitre{Probabilité continue}
\sousChapitre{Théorème central limite}
\difficulte{2}
\theme{théorème central limite, loi exponentielle, approximation}

\contenu{
\begin{question}
\qtype{onlyone}
On note $\mathcal{N}(\mu,\sigma)$ la loi normale d'espérance $\mu$ et d'écart-type $\sigma$, et $\Phi$ la fonction de répartition de $\mathcal{N}(0,1)$. Les durées de traitement de $100$ dossiers, en heures, sont des variables $X_1, \dots, X_{100}$ indépendantes de loi exponentielle de paramètre $1$ (densité $x \mapsto e^{-x}$ sur $[0;+\infty[$, espérance $1$, variance $1$). On note $S = X_1 + \dots + X_{100}$ la durée totale. D'après le théorème central limite, $\mathbb{P}(S \leq 110)$ vaut environ :
\begin{answers}
\good{$\Phi(1)$}
\bad{$\Phi(0{,}1)$}
\bad{$\Phi(10)$}
\bad{$1 - e^{-110}$}
\end{answers}
\begin{explanations}
$\mathbb{E}(S) = 100$ et $\mathbb{V}(S) = 100$, donc
$S$ suit approximativement $\mathcal{N}(100, 10)$ et
$\mathbb{P}(S \leq 110) \approx \Phi\left(\frac{110 - 100}{10}\right) = \Phi(1) \approx 0{,}84$.

B : écart-type $100$, obtenu en multipliant l'écart-type par $n$.

C : écart-type $1$, celui d'une seule durée.

D applique à la somme la fonction de répartition d'une seule durée : $S$ ne suit pas une loi exponentielle.

Contre-sens visés : se tromper d'écart-type pour la somme (B, C) ; croire
que la somme garde la loi de ses termes (D).
\end{explanations}
\end{question}
}
